\chapter{Holomorphic sections and Hermitian mass}
\label{chap:sections}

Let $X$ be a connected compact complex manifold of dimension $n\geq1$, let $L$ be a holomorphic line bundle with smooth Hermitian metric $h$, and fix $u\in\QPSH(X)$. Choose a smooth positive volume form $\mathrm{d}V$.

In a local holomorphic frame $e$, write $|e|_h^2=\mathrm{e}^{-\chi}$ and set $\beta=\ddc\chi$. This defines the smooth closed curvature form of $h$.

The weight selects global sections whose vanishing compensates for its singularities. Their Bergman kernels produce positive currents in the closed curvature class of $L$. A Gauduchon degree bounds how much vanishing a nonzero section can carry; on a submanifold, this gives a necessary condition for realizing the intrinsic trace data by restricted sections.

When $\beta+\ddc u\geq0$ on a curve and $u$ has finitely many logarithmic poles, Riemann--Roch identifies the leading section growth with the residual curvature mass. Higher-dimensional examples explain why the choice of background matters. Homogeneous sections on a Hopf manifold detect its projective base, while an infinite-mass weight with trivial multiplier ideals leaves every section space unchanged. Section counts alone therefore cannot recover a mass computed in an unrelated Hermitian background.

\section{Weighted sections and Bergman currents}\label{sec:weighted-section-spaces}\index{Holomorphic section!weighted space}\label{sec:weighted-bergman-kernel}\index{Bergman kernel!weighted}

The multiplier ideal records which sections have finite weighted norm. Evaluation on the resulting Hilbert space gives the kernel; its extremal formula turns comparisons of norms into comparisons of potentials.

For every positive integer $m$, set
\[
 \mathcal H_m(u)=H^0(X,L^m\otimes\mathcal I(mu)),\qquad
 \|s\|_{m,u}^2=\int_X|s|_{h^m}^2\mathrm{e}^{-mu}\,\mathrm{d}V.
\]
The integral is taken off the pole set of $u$, a set of smooth volume zero. Membership in $\mathcal H_m(u)$ is equivalent to finiteness of this norm: the definition of the multiplier ideal gives local integrability, and compactness gives a finite covering. The converse follows by restriction to coordinate balls. Thus $\mathcal H_m(u)$ is a finite-dimensional Hilbert space, as a subspace of $H^0(X,L^m)$. Write $N_m=\dim\mathcal H_m(u)$.

For $v\in\QPSH(X)$, the ideal order $u\preceq_{\mathcal I}v$ gives
$\mathcal H_m(u)\subseteq\mathcal H_m(v)$ for every $m$. Ordinary equivalence of weights gives the same section spaces. Their Hilbert norms satisfy the following comparison.

\begin{proposition}\label{prop:section-norm-comparison}
Let $v\in\QPSH(X)$, let $D,E\geq0$, and let $F\in C^\infty(X,\mathbb R)$. Suppose $|u-v|\leq D$ off the poles and $\mathrm{d}V'=\mathrm{e}^F \mathrm{d}V$ with $|F|\leq E$. The section spaces for $(u,\mathrm{d}V)$ and $(v,\mathrm{d}V')$ coincide. For every $s$ in this common space and every positive integer $m$,
\[
 \mathrm{e}^{-mD-E}\|s\|_{m,u}^2
 \leq\|s\|_{m,v,\mathrm{d}V'}^2
 \leq \mathrm{e}^{mD+E}\|s\|_{m,u}^2.
\]
\end{proposition}
\begin{proof}
The two integrands differ by the factor $\mathrm{e}^{-m(v-u)+F}$, which lies between the displayed constants. Integrability in either norm is therefore equivalent to integrability in the other. Integrating proves both estimates.
\end{proof}

Fix $m\geq1$ and choose an orthonormal basis $s_1,\ldots,s_{N_m}$ of $\mathcal H_m(u)$. Define $K_{m,u}$ on $X$ and $B_{m,u}$ off the pole set by
\[
 K_{m,u}(x)=\sum_{j=1}^{N_m}|s_j(x)|_{h^m}^2,
 \qquad B_{m,u}(x)=K_{m,u}(x)\mathrm{e}^{-mu(x)}.
\]
Set $B_{m,u}=0$ on the pole set. When $N_m=0$, both kernels are zero. Unitary invariance makes the definition independent of the orthonormal basis.

\begin{proposition}\label{prop:finite-bergman-identities}
If $N_m>0$, then
\[
 K_{m,u}(x)=\sup_{0\neq s\in\mathcal H_m(u)}
                   \frac{|s(x)|_{h^m}^2}{\|s\|_{m,u}^2},
 \qquad \int_X B_{m,u}\,\mathrm{d}V=N_m.
\]
The potential $\Phi_m=m^{-1}\log K_{m,u}$ is genuine, has analytic singularities, and satisfies
\[
 \beta+\ddc\Phi_m\geq0.
\]
This positive current is closed. If $N_m=0$, then $\Phi_m\equiv-\infty$.
\end{proposition}
\begin{proof}
Write a section in the orthonormal basis. The Cauchy--Schwarz inequality bounds its evaluation by the sum of squared basis evaluations, with equality for a coefficient vector proportional to their complex conjugates after trivializing the one-dimensional fiber. At a common zero both sides vanish. This proves the extremal formula. Integrating the defining finite sum for $B_{m,u}$ sums the squared norms of the basis vectors and gives $N_m$.

\Needspace{5\baselineskip}
Locally $s_j=f_j e^{\otimes m}$, so
\[
 \Phi_m=\frac1m\log\sum_j|f_j|^2-\chi,
 \qquad
 \beta+\ddc\Phi_m=\frac1m\ddc\log\sum_j|f_j|^2\geq0.
\]
A nonzero section prevents the sum from vanishing identically. The local expression therefore defines a genuine potential of analytic singularity type. It also proves closedness directly.
\end{proof}

For the weights and volume forms in \cref{prop:section-norm-comparison}, the extremal formula gives, whenever the common section space is nonzero,
\[
 |\Phi_m(u,\mathrm{d}V)-\Phi_m(v,\mathrm{d}V')|\leq D+E/m
\]
off the common base locus. The two potentials have the same ordinary singularity type. In particular the smooth volume form affects these potentials by $O(m^{-1})$ when the weight is fixed.

A simultaneous change of metric and weight preserves the weighted norm exactly. Let $a\in C^\infty(X,\mathbb R)$. If the squared metric is changed to $h'=h \mathrm{e}^{-a}$ and the weight to $u'=u-a$, their product $h'^m \mathrm{e}^{-mu'}$ is unchanged. Hence the Hilbert norm and the density $B_{m,u}\,\mathrm{d}V$ are unchanged, while
\[
 \beta'=\beta+\ddc a,\qquad \Phi'_m=\Phi_m-a.
\]
The current $\beta+\ddc\Phi_m$ is consequently unchanged.

\section{Gauduchon degree and vanishing}\label{sec:gauduchon-section-degree}

The zero divisor of a section is positive. Pairing its Poincar\'e--Lelong formula with a Gauduchon metric therefore bounds its size by the degree of the line bundle. Write $h^0(X,L^m)=\dim H^0(X,L^m)$.

For a Gauduchon metric $\gamma$ on $X$, define the degree
\[
 \deg_\gamma L=\int_X\beta\wedge\gamma^{n-1}.
\]
Changing $h$ changes $\beta$ by $\ddc a$ for a smooth real function $a$. Its integral against $\gamma^{n-1}$ vanishes by Stokes, so the degree is independent of the smooth metric on $L$.

\begin{theorem}\label{thm:gauduchon-section-obstruction}\index{Gauduchon degree}\index{Holomorphic section!vanishing obstruction}
If $\deg_\gamma L<0$, then $H^0(X,L^m)=0$ for every $m\geq1$. If $\deg_\gamma L=0$ and $H^0(X,L^m)\neq0$, then $L^m$ is holomorphically trivial and $h^0(X,L^m)=1$. At degree zero, each nonzero weighted subspace $\mathcal H_m(u)$ is also one-dimensional.
\end{theorem}
\Needspace{7\baselineskip}
\begin{proof}
Let $s$ be a nonzero section of $L^m$ with zero divisor $D$. The Poincar\'e--Lelong formula gives
\[
 \ddc\log|s|_{h^m}^2=[D]-m\beta.
\]
The logarithm is locally integrable. Testing against $\gamma^{n-1}$ gives
\[
 m\deg_\gamma L=\int_D\gamma^{n-1}\geq0.
\]
Every nonempty effective divisor has strictly positive integral because $\gamma$ is positive on its regular locus. Negative degree is therefore impossible. At degree zero the divisor is empty, so $s$ is nowhere zero and trivializes $L^m$. All holomorphic functions on the connected compact manifold are constant, giving dimension one.
\end{proof}

The same degree also constrains positive curvature currents. If $T=\beta+\ddc u\geq0$, the trace identity of \cref{prop:gauduchon-trace} gives
\[
 \deg_\gamma L=\int_X T\wedge\gamma^{n-1}\geq0.
\]
If the degree is zero, positivity and the strict positivity of $\gamma$ imply $T=0$.

\section{Trace data constrains restricted sections}

A restricted section must vanish along the restricted multiplier ideal. Passing to the asymptotic orders of \cref{chap:trace-data} gives a degree obstruction on every smooth compact modification of the submanifold. Subtracting any finite collection of these trace orders also leaves a pseudoeffective class. Here $c_1$ denotes the real Bott--Chern first Chern class.

Let $Y\subset X$ be smooth, compact and connected, of positive dimension $r$, with $\nu(u,Y)=0$. For a positive integer $m$, write
\[
 \mathfrak a_m=\mathcal I(mu)\mathcal O_Y,
 \qquad
 V_m=\operatorname{im}\bigl(\mathcal H_m(u)\longrightarrow H^0(Y,L^m|_Y)\bigr).
\]
Thus $V_m\subseteq H^0(Y,L^m|_Y\otimes\mathfrak a_m)$; the right-hand space can contain sections that do not extend to $X$.

\Needspace{18\baselineskip}
\begin{theorem}\label{thm:gauduchon-trace-section-obstruction}
Suppose $V_m\neq0$ for arbitrarily large integers $m$. Let $\pi:Z\to Y$ be any smooth compact modification, let $E_1,\ldots,E_q$ be distinct prime divisors on $Z$, and let $\gamma$ be a Gauduchon metric on $Z$. Then
\[
 \deg_\gamma(\pi^*(L|_Y))
 \geq\sum_{i=1}^q\tau_{Y,u}(E_i)\int_{E_i}\gamma^{r-1}.
\]
Moreover the real closed $(1,1)$ class
\[
 c_1(\pi^*(L|_Y))-
       \sum_{i=1}^q\tau_{Y,u}(E_i)c_1(\mathcal O_Z(E_i))
\]
is pseudoeffective: it admits a closed positive current.
\end{theorem}
\begin{proof}
Choose a nonzero $s_m\in V_m$ for each integer in an unbounded sequence of such degrees. Its pullback is a nonzero holomorphic section of $\pi^*(L^m|_Y)$. Set $b_i(m)=\operatorname{ord}_{E_i}\mathfrak a_m$, the generic order of the pulled-back image ideal along $E_i$. Since the local coefficient of the pulled-back section belongs to $\mathfrak a_m$, its zero divisor satisfies
\[
 \operatorname{div}(\pi^*s_m)
 \geq\sum_{i=1}^q b_i(m)E_i.
\]
Poincar\'e--Lelong and the Gauduchon identity, exactly as in \cref{thm:gauduchon-section-obstruction}, give
\[
 m\deg_\gamma(\pi^*(L|_Y))
 \geq\sum_{i=1}^q b_i(m)\int_{E_i}\gamma^{r-1}.
\]
Divide by $m$ and use the full integer limit in \cref{thm:intrinsic-trace-profile}. The same limit holds on the selected unbounded sequence, proving the inequality.

Choose smooth metrics on the divisor line bundles and let $\beta_i$ denote their curvatures. The smooth closed form
\[
 \Theta=\pi^*(\beta|_Y)-\sum_{i=1}^q\tau_{Y,u}(E_i)\beta_i
\]
represents the displayed class. Poincar\'e--Lelong shows
$\int_Z\beta_i\wedge\gamma^{r-1}=\int_{E_i}\gamma^{r-1}$.
Thus $\int_Z\Theta\wedge\gamma^{r-1}\geq0$ for every Gauduchon metric on $Z$. For $r\geq2$, the generalized Lamari criterion of \cref{thm:lamari} supplies a genuine $w$ with $\Theta+\ddc w\geq0$. For $r=1$, choose a smooth probability form $\eta$ and solve the scalar Poisson equation $\ddc w=(\int_Z\Theta)\eta-\Theta$; the integral is nonnegative by the same inequality, so the resulting representative is positive. Since $\Theta$ is closed, this is a closed positive current, proving pseudoeffectivity.
\end{proof}

Equivalently, a strict violation of one of these finite Gauduchon inequalities forces $V_m=0$ for all sufficiently large $m$. Indeed the finite sum $\sum_i b_i(m)m^{-1}\int_{E_i}\gamma^{r-1}$ then exceeds the degree for every large $m$, contradicting the inequality satisfied by a nonzero restricted section.

The necessary condition has no general converse. On the elliptic curve $E=\C/(\mathbb Z+\mathrm{i}\mathbb Z)$, choose an irrational real number $a$ and the flat line bundle with multipliers $\mathrm{e}^{2\pi\mathrm{i} a}$ and $1$ on the two generators. Its constant metric has curvature zero. Take $u=0$ and $Y=X=E$, so every trace number is zero and the displayed pseudoeffectivity condition holds. Nevertheless every positive tensor power has no nonzero holomorphic section. To see this, a hypothetical section of a power has a zero divisor of degree zero and hence is nowhere zero. Its logarithmic norm is harmonic on $E$, and is constant. The lifted holomorphic function therefore has constant modulus and is constant on $\C$. The nontrivial multiplier $\mathrm{e}^{2\pi\mathrm{i} ma}$ forces that constant to be zero, a contradiction.


\section{Bounded weighted sections}\label{sec:bounded-section-algebra}

Integrability of two weighted squared norms does not ensure integrability of their product. The logarithmic example below shows how this prevents multiplication of the selected section spaces. Bounded weighted norms are preserved under multiplication and give the graded algebra used in the next chapter.

\begin{example}\label{ex:multiplier-sections-not-graded}\index{Holomorphic section!failure of graded multiplication}
Put $X=\mathbb P^1$, let $z$ be its affine coordinate, and let $\delta_0$ be the Dirac mass at the origin. Take $L=\mathcal O(1)$ with Fubini--Study curvature $\beta$, and let
\[
 u(z)=\frac35\log\frac{|z|^2}{1+|z|^2}.
\]
Then $\beta+\ddc u=\frac25\beta+\frac35\delta_0\geq0$,
$\mathcal I(u)=\mathcal O$ and $\mathcal I(2u)=\mathcal O(-[0])$. A section of $L$ nonzero at $0$ belongs to $\mathcal H_1(u)$, but its square does not belong to $\mathcal H_2(u)$.
\end{example}
\begin{proof}
The weight has only one singularity, $\frac35\log|z|^2+O(1)$ at zero. The one-variable integrability calculation gives the two stated ideals. The curvature follows from Poincar\'e--Lelong and the Fubini--Study potential. The section assertion is precisely the failure to vanish at zero.
\end{proof}

Return to the original data $(X,L,h,u)$. Define $\mathcal W_0=\mathbb C$ and, for $m\geq1$,
\[
 \mathcal W_m(u)=\{s\in H^0(X,L^m):
             |s|_{h^m}^2\mathrm{e}^{-mu}\text{ is bounded off }\{u=-\infty\}\}.
\]
These finite-dimensional vector spaces are unchanged when the smooth metric is changed, since the ratio of two smooth metrics is bounded above and below on $X$. They lie in $\mathcal H_m(u)$, because $X$ has finite smooth volume. For positive integers $m,k$, $s\in\mathcal W_m(u)$ and $t\in\mathcal W_k(u)$, the identity
\[
 |st|_{h^{m+k}}^2\mathrm{e}^{-(m+k)u}
 =\bigl(|s|_{h^m}^2\mathrm{e}^{-mu}\bigr)
  \bigl(|t|_{h^k}^2\mathrm{e}^{-ku}\bigr)
\]
proves $\mathcal W_m(u)\mathcal W_k(u)\subseteq\mathcal W_{m+k}(u)$. Closure under addition follows from $|s+t|^2\leq2|s|^2+2|t|^2$. Thus $\bigoplus_m\mathcal W_m(u)$ is a graded subalgebra of the section ring.

On a smooth projective variety, the almost-semigroup construction of \cite[Section~10.3]{SGPT} treats metrics with positive curvature current on big line bundles with positive multiplier-ideal volume. It uses the integrable section spaces themselves.

\section{Section growth on a curve}

At a logarithmic pole on a curve, integrability rounds the required vanishing order down, while boundedness rounds it up. The resulting divisor degrees differ by a bounded amount. Riemann--Roch therefore gives the same leading growth for both section spaces.

Assume in this section that $X$ is a Riemann surface and $\beta+\ddc u\geq0$. Let $p_1,\ldots,p_r$ be distinct points such that $u$ is locally bounded on $X\setminus\{p_1,\ldots,p_r\}$ and, in coordinates $z_i$ centered at $p_i$,
\[
 u=c_i\log|z_i|^2+O(1),\qquad c_i>0.
\]
Set $\ell=\deg L=\int_X\beta$ and $d=\sum_i c_i$.

\Needspace{14\baselineskip}
\begin{theorem}\label{thm:curve-weighted-section-growth}\index{Holomorphic section!growth on a curve}
In this situation $0\leq d\leq\ell$, and
\[
 \begin{split}
 \mathcal H_m(u)&=H^0\left(X,L^m\Bigl(-\sum_i\lfloor mc_i\rfloor p_i\Bigr)\right),\\
 \mathcal W_m(u)&=H^0\left(X,L^m\Bigl(-\sum_i\lceil mc_i\rceil p_i\Bigr)\right).
 \end{split}
\]
Moreover,
\[
 \lim_{m\to\infty}\frac{\dim\mathcal H_m(u)}m
 =\lim_{m\to\infty}\frac{\dim\mathcal W_m(u)}m
 =\ell-d
 =\int_X\mathrm{MA}^{\mathrm{ext}}_\beta(u).
\]
\end{theorem}
\begin{proof}
A germ vanishing to order $k$ has weighted squared norm comparable to $|z_i|^{2(k-mc_i)}$. Local integrability means $k-mc_i>-1$, or $k\geq\lfloor mc_i\rfloor$; boundedness means $k\geq mc_i$, or $k\geq\lceil mc_i\rceil$. Compactness gives the displayed global identifications.

The positive curvature measure has atom $c_i$ at $p_i$ and total mass $\ell$. Subtracting these atoms leaves a positive measure, proving $d\leq\ell$. The remainder does not charge polar sets: after subtracting the logarithmic singularity, the local potential has a bounded quasi-plurisubharmonic representative. By \cref{prop:bounded-products}, its curvature gives polar sets zero mass. Thus this remaining measure is exactly the extended measure, of mass $\ell-d$.

If $\ell>d$, the degrees of both line bundles in the display equal $m(\ell-d)+O(1)$ and tend to infinity. Let $g$ be the genus of $X$. Riemann--Roch and Serre duality give $h^0=\deg+1-g$ once the degree exceeds $2g-2$. Dividing by $m$ proves both limits.

If $\ell=d$, both degrees remain bounded. A line bundle of negative degree has no nonzero section, and one of nonnegative degree $k$ has at most $k+1$ sections: impose vanishing at $k+1$ distinct points and use the negative-degree case for the kernel. Hence both dimensions are uniformly bounded, and both limits are zero. These arguments include $r=0$.
\end{proof}

\Needspace{9\baselineskip}
\section{Homogeneous sections on a Hopf manifold}\label{sec:hopf-homogeneous-sections}

On a scalar Hopf manifold, a holomorphic section of a line bundle lifts to a function on $\mathbb C^n\setminus\{0\}$ with an exact homogeneity property under the contraction. Hartogs extension reduces the section spaces to spaces of homogeneous polynomials, and the invariance of the metric makes the Bergman kernel constant. The resulting growth is of order $m^{n-1}$, reflecting the dimension of the projective base rather than that of the manifold.

Let $n\geq2$, let $z=(z_1,\ldots,z_n)$ be the coordinates on $\mathbb C^n$, and set
\[
 X_n=(\C^n\setminus\{0\})/\langle z\mapsto \mathrm{e}^{-1}z\rangle.
\]

\Needspace{6\baselineskip}
On the covering space put
\[
 s=\log|z|^2,\qquad \alpha=\ddc s,\qquad
 \zeta=\mathrm{d}s\wedge\dc s,\qquad \omega=\alpha+\zeta.
\]
The contraction changes $s$ by the constant $-2$, so these forms descend to $X_n$. The calculation proving \eqref{eq:hopf-torsion}, written in \cref{chap:hopf} for $n=3$ and with $\beta$ in place of $\zeta$, applies in every dimension and gives
\[
 \alpha^n=0,\qquad \zeta^2=0,\qquad
 \ddc\zeta=-\alpha^2,
 \qquad \ddc(\omega^{n-1})=-(n-1)\alpha^n=0.
\]
Thus $\omega$ is a Gauduchon form. Projectivization defines a holomorphic map $\pi\colon X_n\to\mathbb P^{n-1}$ with elliptic fibers, and $\alpha$ is the pullback of the Fubini--Study form normalized to have total mass one. Since $\zeta$ restricts to twice the normalized area form on each fiber, as in the proof of \cref{prop:hopf-mass}, and $\omega^n=n\alpha^{n-1}\wedge\zeta$, we obtain
\[
 \int_{X_n}\alpha^{n-1}\wedge\zeta=2,
 \qquad \int_{X_n}\omega^n=2n.
\]
Set $\mathrm{d}V=\omega^n/(2n)$, a probability measure.

For an integer $k\geq0$, let $L_k$ be the line bundle on $X_n$ obtained from $(\C^n\setminus\{0\})\times\C$ through the identification
\[
 (z,t)\sim(\mathrm{e}^{-1}z,\mathrm{e}^{-k}t).
\]
Give it the metric $|t|_{h_k}^2=|t|^2|z|^{-2k}$. Under the identification the two factors change by $\mathrm{e}^{-2k}$ and $\mathrm{e}^{2k}$, so this metric descends. In the frame $t=1$ the squared norm is $|z|^{-2k}$, and the curvature form is $k\alpha$. We apply the constructions of \cref{sec:weighted-section-spaces} with $X=X_n$, $L=L_k$, $h=h_k$ and this probability volume form.

\Needspace{15\baselineskip}
\begin{proposition}\label{prop:hopf-bergman-exact}\index{Bergman kernel!homogeneous Hopf example}
For $m\geq1$ and the weight $u=0$,
\[
 N_m:=h^0(X_n,L_k^m)=\binom{km+n-1}{n-1},
 \qquad B_{m,0}=N_m.
\]
The Bergman potential is the constant $m^{-1}\log N_m$, and the Bergman current is $k\alpha$. Moreover,
\[
 \deg_\omega L_k=2k(n-1),\qquad
 \lim_{m\to\infty}m^{-n}N_m=0,
 \qquad (k\alpha)^n=0,
\]
whereas $\int_{X_n}\omega^n=2n>0$.
\end{proposition}
\begin{proof}
A section of $L_k^m$ lifts to a holomorphic function $f$ on $\C^n\setminus\{0\}$ satisfying
\[
 f(\mathrm{e}^{-1}z)=\mathrm{e}^{-km}f(z).
\]
Since $n\geq2$, Hartogs' theorem extends $f$ holomorphically to $\C^n$. In its Taylor series, a monomial of total degree $j$ is multiplied by $\mathrm{e}^{-j}$ under the contraction, so comparing coefficients shows that only monomials of degree $km$ occur. Conversely, every homogeneous polynomial of that degree has the required equivariance. Counting monomials gives the binomial coefficient; when $k=0$, only the constant polynomials remain.

The inner product defined by $h_k^m$ and $\mathrm{d}V$ is invariant under the unitary group $U(n)$, so the basis-independent function $K_{m,0}$ is $U(n)$-invariant. It is also invariant under positive radial dilations: every section is homogeneous of degree $km$, and the factor $|z|^{-2km}$ in its squared norm cancels the effect of a dilation. Unitary transformations and positive radial dilations act transitively on $\C^n\setminus\{0\}$, so $K_{m,0}$ is constant. By \cref{prop:finite-bergman-identities} its integral is $N_m$, and $\mathrm{d}V$ has mass one. Thus $K_{m,0}=B_{m,0}=N_m$.

\Needspace{6\baselineskip}
The statements about the Bergman potential and current follow immediately. Since $\alpha^n=0$ and $\zeta^2=0$,
\[
 \int_{X_n}k\alpha\wedge\omega^{n-1}
 =k(n-1)\int_{X_n}\alpha^{n-1}\wedge\zeta
 =2k(n-1).
\]
The binomial coefficient grows at most like $m^{n-1}$, which proves the normalized limit, and $\alpha^n=0$ proves the vanishing of the curvature mass.
\end{proof}

For $n=3$, $X_3$ is the Gauduchon threefold of \cref{thm:hopf-unbounded}, on which Hermitian Monge--Amp\`ere masses are unbounded; for $k>0$, the line bundles $L_k$ nevertheless have positive degree and quadratic section growth. In the notation of \cref{ex:hopf-exact-complement-index}, $L_k$ is the line bundle $E_a$ on $X_q$ with $q=\mathrm e^{-1}$ and $a=\mathrm e^{-k}$. Hence, for every $n\geq2$ and $k>0$, the complement $\mathbb P(\mathcal O\oplus L_k)\setminus\mathbb P(\mathcal O)$ of \cref{subsec:flat-restriction-hopf-complements} has least smooth convexity index two, that is, it is $2$-convex but not $1$-convex in the sense of \cref{subsec:concave-neighbourhood-dual-norms}. When $k=0$, the line bundle is trivial and the least smooth convexity index of its complement is $n+1$, by the same example.


\section{Section counts and infinite Hermitian mass}

Return to the initial compact setting. A weight with trivial multiplier ideals imposes no integrability restriction on global sections. The infinite-mass witness of \cref{thm:infinite-mass-witness} therefore leaves every section count unchanged. For the trivial bundle, all Bergman currents vanish while the Hermitian mass is infinite.

\Needspace{10\baselineskip}
\begin{proposition}\label{prop:sections-infinite-mass-obstruction}
Let $\omega$ be a Hermitian metric on $X$ with $v_+(\omega)=+\infty$. There is a potential $U\in\PSH(X,\omega)$ such that, with $L=\mathcal O_X$ and its constant metric,
\[
 \dim\mathcal H_m(U)=1\quad\hbox{for every }m\geq1,
 \qquad \int_X\mathrm{MA}^{\mathrm{ext}}_\omega(U)=+\infty.
\]
The currents produced by its global Bergman potentials are identically zero.
\end{proposition}
\begin{proof}
Take $U$ from \cref{thm:infinite-mass-witness}, which has $\mathcal I(tU)=\mathcal O_X$ at every positive weight and infinite extended mass. Global holomorphic functions on $X$ are constant, so every weighted section space is one-dimensional. The constant section is integrable by the ideal condition and compactness. Its normalized basis consists of a nonzero constant; therefore each $\Phi_m$ is constant. Since the trivial metric has curvature zero, its Bergman current is zero.
\end{proof}

In particular $m^{-n}\dim\mathcal H_m(U)\to0$ while the Hermitian mass is infinite. Here the line-bundle curvature is zero and the nonlinear operator uses $\omega$; the two quantities arise from different backgrounds.

The bounded section algebra now supplies the multiplicative structure for the flag bodies of \cref{chap:flag-bodies}. Quantitative kernel limits require further relations between the selecting weight and line-bundle curvature; these are developed in \cref{chap:partial-bergman}.
